\section{Conclusion}

\myindent In this internship project under the supervision of Michele Castellana and David Lacoste, we implemented a bitwise version of the Metropolis algorithm. This version uses the machine-specific binary arithmetic to efficiently simulate the evolution of $n_\text{bits}$ systems in parallel. On top of this structural optimization, the same sequence of random numbers is used to make all $n_\text{bits}$ systems evolve. This turns out to be central as random number generation is often the simulation bottleneck. Using this bitwise formalism considerably reduces the number of operations required to make the system evolve, as all configurations are stored using \Bits objects, whose evolutions are given by bitwise operations that are straightforward for the machine. For now, we use $n_\text{bits}=64$, but this number can potentially be increased up $n_\text{bits}=512$, providing a change in the type of objects used and low-level implementation of the code. The limitations of the size of $n_\text{bits}$ are directly related to the constraints of the CPU architecture.\\
\myindent For now, the code only provides a bitwise implementation of the Metropolis algorithm, but other Monte-Carlo algorithm such as the Glauber algorithm \cite{Glauber1963} or cluster flip algorithms \cite{SwendsenWang1987,Wolff1989} can also be setup in this formalism. So far, the code enables the application of this bitwise Metropolis algorithm on several networks models: the Ising model, the Edwards-Anderson spin glass model and the Hopfield model. It can however be generalized to every model that uses discrete values, as it is always possible to shift those values into positive integers that are suitably stored in \UnsignedInt objects.\\
\myindent The interest of this bitwise formalism is to efficiently sample the dynamics of a given system. Indeed, the Metropolis algorithm becomes computationally demanding when one attempts to sample a large number of configurations. This situation occurs, for example, when looking for rare samples in the dynamics, or sampling a large number of initial conditions.\\
\myindent This implementation therefore seemed ideal for studying the dynamics of the Hopfield model. Indeed, while the Hopfield model -- introduced in its now-canonical form by Hopfield's landmark 1982 paper \cite{Hopfield1982} -- has become a reference in the statistical mechanics of networks and a milestone in research on both biological and artificial neural networks, its dynamical properties remain for the most part to be discovered. The objective of this internship was thus to use this implementation to study the dynamics of the model.\\
\myindent The code has first been tested on the Ising and Edwards-Anderson models, for which it reproduced the correct expected results, see \Cref{fig:phase_diagram_ising} and \Cref{fig:binder}. In addition, it has proven to be efficient, yielding computational gains -- that is the ratio of the number of independent results obtained for a given computation time between the optimized bitwise implementation and the reference one -- comprised between 20 and 45, depending mostly on the temperature of the system. However, when applying the same bitwise implementation on the Hopfield network, the computational gain vanished. This is most likely due to the density of the network: one has to go through every single neuron in the network to update a single neuron. If the bitwise implementation creates some temporary \Bits or \UnsignedInt objects during the spin interaction calculations, the loss in performance can become considerable for dense networks. Fixing this low-level implementation in order to prevent the creation of temporary objects will be the focus of the remaining two weeks of this internship. Despite this loss of computational advantage, the bitwise implementation still yielded the expected results, see \Cref{fig:reconstruction_error}. In addition, this implementation allows us to easily visualize the evolution of the $n_\text{bits}$ independent system in parallel, see \Cref{fig:evolution_ising}, \Cref{fig:all_mu_single_replica} and \Cref{fig:mu_star_all_replica}.\\
\myindent This method thus proved promising to study the dynamics of the Hopfield network, that is the dynamical exponent, correlations or even the dependence of this dynamics on the network's structure. Unfortunately, this project will most likely not be pursued into a PhD project, due to the lack of fundings.
